%auto-ignore
\section{절제 연구}
\label{sec:ablation}
본 절에서 저자들은 각 요소의 상대적 중요도를 더 잘 이해하기 위해 BERT의 여러 측면에 대한 절제 실험(ablation study)을 수행한다. 추가 절제 연구는 부록~\ref{appendix:sec:more_ablation_studies}에서 확인할 수 있다.

\subsection{사전학습 과제의 효과}

\label{sec:task_ablation}
저자들은 \bertbase와 정확히 동일한 사전학습 데이터, 파인튜닝 방식, 하이퍼파라미터를 쓰면서 두 가지 사전학습 목적함수를 평가해 BERT의 깊은 양방향성의 중요성을 입증한다.
\vspace{0.3cm}
\\
\noindent\textbf{No NSP}: ``마스킹 LM''(MLM)으로 학습되지만 ``다음 문장 예측''(NSP) 과제 없이 학습되는 양방향 모델.\\
\noindent\textbf{LTR \& No NSP}: MLM이 아닌 표준 좌→우(LTR) LM으로 학습되는 좌측 맥락만 사용하는 모델. 좌측 전용 제약은 파인튜닝에도 동일하게 적용했는데, 이 제약을 제거하면 사전학습/파인튜닝 불일치가 생겨 다운스트림 성능이 떨어졌기 때문이다. 또한 이 모델은 NSP 과제 없이 사전학습된다. 이 모델은 OpenAI GPT와 직접 비교 가능하지만, 더 큰 학습 데이터셋과 본 논문의 입력 표현·파인튜닝 방식을 사용한다.
%auto-ignore
\begin{table}[t]
\small
 \begin{tabular}{@{}lccccc@{}}
    \toprule
              & \multicolumn{5}{c}{Dev 셋} \\
   과제 & MNLI-m & QNLI & MRPC & SST-2 & SQuAD     \\
         & (Acc) & (Acc) & (Acc) & (Acc) & (F1)     \\
     \midrule
\bertbase       & 84.4 & 88.4 & 86.7 & 92.7 & 88.5 \\
No NSP          & 83.9 & 84.9 & 86.5 & 92.6 & 87.9 \\
LTR \& No NSP   & 82.1 & 84.3 & 77.5 & 92.1 & 77.8 \\
\quad + BiLSTM  & 82.1 & 84.1 & 75.7 & 91.6 & 84.9 \\
     \bottomrule
   \end{tabular}
   \caption{\bertbase 아키텍처를 사용한 사전학습 과제 절제. ``No NSP''는 다음 문장 예측 과제 없이 학습된다. ``LTR \& No NSP''는 OpenAI GPT처럼 다음 문장 예측 없이 좌→우 LM으로 학습된다. ``+ BiLSTM''은 파인튜닝 시 ``LTR + No NSP'' 모델 위에 무작위 초기화된 BiLSTM을 추가한 것이다.
   }
   \label{tab:task_ablation}    
\end{table}

먼저 NSP 과제의 영향을 살핀다. 표~\ref{tab:task_ablation}에서 보이듯, NSP를 제거하면 QNLI, MNLI, SQuAD 1.1에서 성능이 크게 떨어진다. 다음으로, ``No NSP''와 ``LTR \& No NSP''를 비교해 양방향 표현 학습의 영향을 평가한다. LTR 모델은 모든 과제에서 MLM 모델보다 성능이 낮으며, MRPC와 SQuAD에서 큰 폭의 하락을 보인다.

SQuAD에서는 LTR 모델이 토큰 예측에서 좋지 않은 성능을 낼 것임이 직관적으로 분명하다 — 토큰 수준 은닉 상태에 우측 맥락이 전혀 없기 때문이다.
LTR 시스템을 강화하려는 성실한 시도로, 저자들은 그 위에 무작위로 초기화된(randomly initialized) BiLSTM을 추가했다. 이는 SQuAD 결과를 의미 있게 개선하지만, 사전학습된 양방향 모델의 결과에는 한참 못 미친다. BiLSTM은 GLUE 과제에서는 오히려 성능을 떨어뜨린다.

저자들은 ELMo처럼 별도의 LTR과 RTL 모델을 학습한 뒤 각 토큰을 두 모델의 결합으로 표현하는 것도 가능함을 인정한다. 그러나 (a) 이는 단일 양방향 모델보다 두 배 비싸고, (b) RTL 모델이 답을 질문에 조건부로 둘 수 없으므로 QA 같은 과제에는 비직관적이며, (c) 모든 층에서 좌·우 맥락을 모두 사용할 수 있는 깊은 양방향 모델보다 엄격하게 약하다.


\subsection{모델 크기의 효과}
\label{sec:model_size_ablation}

본 절에서 저자들은 모델 크기가 파인튜닝 과제 정확도에 미치는 영향을 탐구한다. 층수, 은닉 유닛 수, 어텐션 헤드 수가 다른 여러 BERT 모델을 학습했고, 그 외의 하이퍼파라미터와 학습 절차는 앞서 기술한 것과 동일하게 두었다.

선택된 GLUE 과제의 결과는 표~\ref{tab:size_ablation}에 있다. 본 표에서 저자들은 5번의 무작위 재시작 파인튜닝의 평균 Dev 셋 정확도를 보고한다. 더 큰 모델이 네 데이터셋 모두에서 엄격한 정확도 향상을 가져옴을 알 수 있다. 학습 라벨이 3,600개에 불과하고 사전학습 과제와 상당히 다른 MRPC조차도 마찬가지다. 또한 기존 문헌에 비해 이미 상당히 큰 모델 위에서 이 정도 향상을 얻을 수 있다는 점은 다소 놀라울 수 있다. 예컨대 \citet{vaswani-etal:2017:_atten}에서 탐구한 가장 큰 트랜스포머는 인코더에 100M 파라미터를 가지는 ($L=6$, $H=1024$, $A=16$)이고, 문헌에서 저자들이 발견한 가장 큰 트랜스포머는 235M 파라미터의 ($L=64$, $H=512$, $A=2$)이다~\cite{alrfou:2018}. 이에 비해 \bertbase는 110M, \bertlarge는 340M 파라미터를 갖는다.

%auto-ignore
\begin{table}[b]
\begin{center}
{\small
\begin{tabular}{@{}rrrcccc@{}}
  \toprule
  \multicolumn{3}{c}{하이퍼파라미터}      &      & \multicolumn{3}{c}{Dev 셋 정확도} \\
  \midrule
  \#L & \#H &\#A & LM (ppl) & MNLI-m & MRPC &SST-2               \\
  \midrule
  \
   3 &  768 & 12 & 5.84 & 77.9 & 79.8 & 88.4 \\
   6 &  768 &  3 & 5.24 & 80.6 & 82.2 & 90.7 \\
   6 &  768 & 12 & 4.68 & 81.9 & 84.8 & 91.3 \\
  12 &  768 & 12 & 3.99 & 84.4 & 86.7 & 92.9 \\
  12 & 1024 & 16 & 3.54 & 85.7 & 86.9 & 93.3 \\
  24 & 1024 & 16 & 3.23 & 86.6 & 87.8 & 93.7 \\
\bottomrule
\end{tabular}
} % small
\end{center}
\caption{\label{tab:size_ablation} BERT 모델 크기 절제. \#L = 층수, \#H = 은닉 크기, \#A = 어텐션 헤드 수. ``LM (ppl)''은 홀드아웃 학습 데이터의 마스킹 LM 퍼플렉시티다.}
\end{table}


기계 번역과 언어 모델링 같은 대규모 과제에서 모델 크기를 키우면 지속적인 향상이 따라온다는 점은 오래 전부터 알려져 있고, 표~\ref{tab:size_ablation}이 보여 주는 홀드아웃(held-out) 학습 데이터의 LM 퍼플렉시티(perplexity)에서도 이를 확인할 수 있다. 그러나 모델이 충분히 사전학습되었다는 전제 하에 \emph{매우 작은 규모의 과제}에서도 극단적인 모델 크기로의 스케일링이 큰 향상을 가져온다는 점을 설득력 있게 보인 것은 본 연구가 처음이라고 저자들은 본다. \citet{peters2018dissecting}은 사전학습 bi-LM의 층수를 2에서 4로 늘렸을 때 다운스트림 영향에 대해 혼합된 결과를 제시했고, \citet{melamud2016context2vec}은 은닉 차원을 200에서 600으로 늘리면 도움이 되지만 1,000으로 더 늘려도 추가 향상은 없었다고 짤막하게 언급했다. 이 두 선행 연구는 모두 피처 기반 접근을 사용했다 --- 저자들의 가설은, 모델이 다운스트림 과제 위에서 직접 파인튜닝되고 무작위 초기화된 추가 파라미터가 매우 적게만 도입될 때에는, 다운스트림 데이터가 매우 작더라도 과제별 모델이 더 크고 표현력 있는 사전학습 표현으로부터 이득을 볼 수 있다는 것이다.



\subsection{BERT의 피처 기반 접근}
\label{sec:ner}
지금까지 제시한 BERT 결과는 모두, 사전학습 모델 위에 단순한 분류 층을 더한 뒤 다운스트림 과제 위에서 모든 파라미터를 함께 파인튜닝하는 파인튜닝 접근을 사용했다. 그러나 사전학습 모델에서 고정 피처를 추출하는 피처 기반 접근에도 일정한 장점이 있다. 첫째, 모든
%NLP
과제가 트랜스포머 인코더 아키텍처로 손쉽게 표현되지는 않으며, 따라서 과제별 모델 아키텍처가 추가로 필요할 수 있다. 둘째, 학습 데이터의 비싼 표현을 한 번만 미리 계산해 두고 그 위에서 더 저렴한
%being able to
%less expensive
모델로 여러 실험을 돌릴 수 있다는 큰 계산 이점이 있다.

본 절에서 저자들은 CoNLL-2003 개체명 인식(named entity recognition, NER) 과제~\cite{tjong-de:2003}에 BERT를 적용해 두 접근을 비교한다. BERT 입력에는 대소문자를 보존하는 WordPiece 모델을 사용하며, 데이터가 제공하는 최대 문서 맥락을 함께 포함한다. 표준 관행을 따라 이를 태깅(tagging) 과제로 정식화하지만, 출력에는 CRF(Conditional Random Field) 층을 사용하지 않는다. 첫 서브토큰의 표현을 NER 라벨 집합 위 토큰 수준 분류기의 입력으로 사용한다.


파인튜닝 접근을 절제하기 위해, 저자들은 BERT의 어떤 파라미터도 파인튜닝하지 \emph{않은} 채 하나 이상의 층의 활성을 추출하는 피처 기반 접근을 적용한다. 이렇게 얻은 맥락 임베딩은, 분류 층 이전에 무작위 초기화된 두 층 768차원 BiLSTM의 입력으로 사용된다.

결과는 표~\ref{tab:pretrained_embeddings}에 제시한다. \bertlarge는 최첨단 방법들과 경쟁력 있는 성능을 보인다. 가장 좋은 방법은 사전학습 트랜스포머의 상위 네 개 은닉 층의 토큰 표현을 결합하는 것이며, 이는 모델 전체를 파인튜닝하는 결과보다 0.3 F1만 뒤진다. 이는 BERT가 파인튜닝 접근과 피처 기반 접근 양쪽 모두에 효과적임을 보여 준다.


%auto-ignore

\begin{table}[t]
\small
\centering
 \begin{tabular}{@{}lcc@{}}
\toprule
System & Dev F1 & Test F1 \\
\midrule
ELMo~\cite{peters-etal:2018:_deep}& 95.7 & 92.2 \\
CVT~\cite{clark2018semi} & - & 92.6 \\
CSE~\cite{akbik2018contextual} & - & {\bf 93.1} \\
\midrule
파인튜닝 접근 & & \\
\;\;\;\bertlarge  & 96.6 & 92.8 \\
\;\;\;\bertbase& 96.4 & 92.4 \\
\midrule
피처 기반 접근 (\bertbase) &  &  \\
\;\;\;Embeddings & 91.0 &- \\
\;\;\;Second-to-Last Hidden   & 95.6&- \\
\;\;\;Last Hidden            & 94.9&- \\
\;\;\;Weighted Sum Last Four Hidden        & 95.9&- \\
\;\;\;Concat Last Four Hidden        & 96.1&- \\
\;\;\;Weighted Sum All 12 Layers        & 95.5&- \\
\bottomrule
\end{tabular}
\caption{CoNLL-2003 개체명 인식 결과. 하이퍼파라미터는 Dev 셋을 이용해 선택했다. 보고된 Dev·Test 점수는 그 하이퍼파라미터로 5번 무작위 재시작한 결과의 평균이다.
}
\label{tab:ner_results}    
\label{tab:pretrained_embeddings}    
\end{table}
